\documentclass[10pt,a4paper]{amsart}
\usepackage[margin=19mm]{geometry}
\usepackage{amsmath,amssymb,mathtools}
\usepackage[scr=rsfs,cal=euler]{mathalfa}
\usepackage{xfrac}
\usepackage[unicode,hidelinks,bookmarks=false]{hyperref}
\newcommand{\ie}{i.e.\ }
\newcommand{\cf}{cf.\ }
\newcommand{\resp}{respectively\ }
\newcommand{\Subsection}{\subsection}
\providecommand{\nobreakdash}{\nobreak\hbox{-}}
\setlength{\emergencystretch}{2em}
\multlinegap0pt
\usepackage{bm}
\usepackage{bbm}
\usepackage{dsfont}
\usepackage{xspace}
\usepackage{cancel}
\usepackage{mathtools}
\usepackage{amsthm}
\makeatletter
\@ifundefined{thm}{\newtheorem{thm}{Thm}}{}
\@ifundefined{prop}{\newtheorem{prop}[thm]{Proposition}}{}
\makeatother
\title[Infinitesimally rigid Lie foliations with dense leaves]{Infinitesimally rigid Lie foliations with dense leaves}
\author{Stephane Geudens, Florian Zeiser}
\date{Source: arXiv:2403.17666v6 (CC BY 4.0)}
\begin{document}
\maketitle

\noindent\emph{Source and licence.} Infinitesimally rigid Lie foliations with dense leaves. Version arXiv:2403.17666v6 (CC BY 4.0), distributed under the Creative Commons Attribution 4.0 International licence, as recorded in the arXiv licence metadata for this exact version and shown on the abstract page https://arxiv.org/abs/2403.17666v6. Full rights notice: licenses/arxiv-2403.17666-CC-BY-4.0.txt.\par\smallskip

\noindent\textit{Editorial scope.} Counted unit: the Prop whose source label is \texttt{prop:spectral} in the article (source0.tex lines 751--757, source label \texttt{prop:spectral}), delivered together with the article's own complete original proof of it (source0.tex lines 758--805, one \texttt{proof} environment of 48 lines). No mathematics is written, repaired or re-derived: statement and proof are the article's own text line for line. Required context retained uncounted: eq:E\_1 (lines 724--726); eq:E\_1p (lines 728--730); eq:bas-coho (lines 747--749). Every definition, display and label of those lines is retained unchanged. Omitted from this package and not redistributed: 1--723, 727--727, 731--746, 750--750, 806--1842. Nothing inside a retained excerpt is omitted or altered: every retained body is delivered exactly as the article's lines are written, so no omission receipt of any kind accompanies this package. Citations: the retained excerpts carry no citation command at all, so no bibliography entry of the article is transcribed and no reference is redistributed. Reference adaptation: none was needed and none was made; every reference inside the delivered unit resolves inside it, so no unresolved reference or citation is produced. Printed slips kept literal: no printed word, symbol, hypothesis or number of the counted statement or of its proof was altered. Layout and class: the article's own class is \texttt{amsart} with packages including latexsym, amsmath, amssymb, amsfonts, amscd, bm, amsthm, t1enc, eucal, indentfirst, graphicx, pb-diagram, fancyhdr, fancybox; the wrapper is the lane's pinned \texttt{amsart} editorial wrapper at 10pt on A4, which restates the theorem-like environments the retained text needs and adds the article's own macro definitions that the retained text needs (none), with extra wrapper packages (bm, bbm, dsfont, xspace, cancel, mathtools). The delivered numbering is the unit's own. Assets: no retained span contains a figure environment, a graphics-inclusion command, a PDF-page inclusion, a table or a TikZ picture; no image file is copied into this package, the package declares no asset dependency and no image file is redistributed with it. Licence: version arXiv:2403.17666v6 of this article is distributed under the Creative Commons Attribution 4.0 International licence, as recorded in the arXiv licence metadata for this exact version and shown on the abstract page https://arxiv.org/abs/2403.17666v6. A full rights notice is delivered at licenses/arxiv-2403.17666-CC-BY-4.0.txt. Characters: every retained excerpt is pure ASCII, so the delivered engine, XeLaTeX with its default Latin Modern text font, reports no missing character.
	\begin{equation}\label{eq:E_1}
		E_1^{p,q}\cong H^{q}\big(\mathcal{F},\wedge^{p}(T\mathcal{G}/T\mathcal{F})^{*}\big).
	\end{equation}

	\begin{equation}\label{eq:E_1p}
		E_1^{p,0}\cong\big\{\xi\in\Gamma\big(\wedge^{p}(T\mathcal{G}/T\mathcal{F})^{*}\big):\nabla_X^{*}\xi=0\ \ \forall X\in\Gamma(T\mathcal{F})\big\}.
	\end{equation}

	\begin{equation}\label{eq:bas-coho}
		E_2^{p,0}\cong H^{p}_{\mathcal{F}}(\mathcal{G}).
	\end{equation}

	\begin{prop}\label{prop:spectral}
		Let $\mathcal{F}$ be a Lie $\mathfrak{g}$-foliation on a manifold $M$ defined by a Maurer-Cartan form $\omega\in\Omega^{1}(M,\mathfrak{g})$. Let $\mathfrak{h}\subset\mathfrak{g}$ be an ideal and denote by $pr:\mathfrak{g}\rightarrow \mathfrak{g}/\mathfrak{h}$ the projection. 
		\begin{enumerate}
			\item The form $pr\circ\omega\in\Omega^{1}(M,\mathfrak{g}/\mathfrak{h})$ defines a Lie $\mathfrak{g}/\mathfrak{h}$-foliation $\mathcal{G}$.
			\item Assume that $\mathcal{F}$ is infinitesimally rigid with dense leaves. Then $H^{1}(\mathcal{G})\cong H^{1}(\mathfrak{h})$.
		\end{enumerate}
	\end{prop}

	\begin{proof}
		$(1)$ For every point $x\in M$ we have a surjective map
		\[
		pr\circ\omega_x:T_{x}M\twoheadrightarrow\mathfrak{g}\twoheadrightarrow\mathfrak{g}/\mathfrak{h}.
		\]
		Moreover, $pr\circ\omega\in\Omega^{1}(M,\mathfrak{g}/\mathfrak{h})$ satisfies the Maurer-Cartan equation because the projection $pr:\mathfrak{g}\rightarrow \mathfrak{g}/\mathfrak{h}$ is a Lie algebra homomorphism, i.e.
		\[
		0=pr\circ\left(d\omega+\frac{1}{2}[\omega,\omega]\right)=d\big(pr\circ\omega\big)+\frac{1}{2}[pr\circ\omega,pr\circ\omega].
		\]
		$(2)$ We start by examining the Serre spectral sequence $E_k$ of the Lie subalgebroid $T\mathcal{F}\subset T\mathcal{G}$. Since it converges to $H^{\bullet}(\mathcal{G})$, we have that
		\[
		H^{1}(\mathcal{G})\cong E_{\infty}^{1,0}\oplus E_{\infty}^{0,1}.
		\]
		By the isomorphism \eqref{eq:E_1}, we have that $E_1^{0,1}\cong H^{1}(\mathcal{F})$, which vanishes since $\mathcal{F}$ is assumed to be infinitesimally rigid. Hence, also the limiting term $E_{\infty}^{0,1}$ vanishes. Because the differential $d_k$ of page $E_k$ has bi-degree $(k,1-k)$, it follows that
		\[
		E_{\infty}^{1,0}=E_2^{1,0}\cong H^{1}_{\mathcal{F}}(\mathcal{G}),
		\]
		where we used the isomorphism \eqref{eq:bas-coho}. So the result follows if we show that $H^{1}_{\mathcal{F}}(\mathcal{G})\cong H^{1}(\mathfrak{h})$.
		
		We will show that the complex  $\big(\Omega^{\bullet}_{\mathcal{F}}(\mathcal{G}),d_{\mathcal{G}}\big)$ is isomorphic with the Chevalley-Eilenberg complex $\big(\wedge^{\bullet}\mathfrak{h}^{*},d_{CE}\big)$.
		To do so, recall that the space of transverse fields $\mathfrak{X}(M,\mathcal{F})/\Gamma(T\mathcal{F})$ is canonically a Lie algebra, where $\mathfrak{X}(M,\mathcal{F})$ is the space of $\mathcal{F}$-projectable vector fields. Since the leaves of $\mathcal{F}$ are dense, we have a natural Lie algebra isomorphism
		\[
		\Psi:\big(\mathfrak{g},[\cdot,\cdot]\big)\overset{\sim}{\longrightarrow}\left(\mathfrak{X}(M,\mathcal{F})/\Gamma(T\mathcal{F}),[\cdot,\cdot]\right):v\mapsto \overline{Y_v},
		\]
		where $\overline{Y_v}$ is the unique transverse field satisfying $\omega(\overline{Y_v})=v$.
		Fix a basis $\{v_1,\ldots,v_k\}$ of $\mathfrak{h}$ and let $\overline{Y_1},\ldots,\overline{Y_k}\in\mathfrak{X}(M,\mathcal{F})/\Gamma(T\mathcal{F})$ denote the corresponding transverse fields.
		It follows that the space $\mathbb{R}\overline{Y_1}\oplus\cdots\oplus\mathbb{R}\overline{Y_k}$ inherits a Lie algebra structure and that we have an isomorphism of Lie algebras
		\begin{equation}\label{eq:iso-Lie}
			\Psi:\big(\mathfrak{h},[\cdot,\cdot]\big)\overset{\sim}{\longrightarrow}\left(\mathbb{R}\overline{Y_1}\oplus\cdots\oplus\mathbb{R}\overline{Y_k},[\cdot,\cdot]\right).
		\end{equation}
		Note that any representative $Y_i$ of the class $\overline{Y_i}$ is in fact tangent to $\mathcal{G}$ because
		\[
		(pr\circ\omega)(Y_i)=pr(v_i)=0.
		\]
		It follows that $\{\overline{Y_1},\ldots,\overline{Y_k}\}$ is a frame of $T\mathcal{G}/T\mathcal{F}$. Moreover, each  $\overline{Y_i}$ is flat with respect to the connection $\nabla$ because any representative $Y_i$ of $\overline{Y_i}$ is $\mathcal{F}$-projectable. This implies that the dual frame $\{\xi_1,\ldots,\xi_k\}$ of $(T\mathcal{G}/T\mathcal{F})^{*}$ consists of sections that are flat for $\nabla^{*}$. Because the leaves of $\mathcal{F}$ are dense, the description \eqref{eq:E_1p} of $\Omega^{p}_{\mathcal{F}}(\mathcal{G})$ shows that
		\[
		\Omega^{p}_{\mathcal{F}}(\mathcal{G})\cong\bigoplus_{1\leq i_1<\cdots<i_p\leq k}\mathbb{R}~\xi_{i_1}\wedge\cdots\wedge\xi_{i_p}.
		\]
		Since 
		\[
		d_{\mathcal{G}}\xi_\alpha\big(\overline{Y_\beta},\overline{Y_\gamma}\big)=-\xi_{\alpha}[\overline{Y_\beta},\overline{Y_\gamma}],
		\]
		we obtain that the complex $\big(\Omega^{\bullet}_{\mathcal{F}}(\mathcal{G}),d_{\mathcal{G}}\big)$ is nothing but the Chevalley-Eilenberg complex of the Lie algebra $\left(\mathbb{R}\overline{Y_1}\oplus\cdots\oplus\mathbb{R}\overline{Y_k},[\cdot,\cdot]\right)$. Since the map \eqref{eq:iso-Lie} is a Lie algebra isomorphism, exterior powers of its dual give an isomorphism of Chevalley-Eilenberg complexes
		\[
		\wedge^{\bullet}\Psi^{*}:\big(\Omega^{\bullet}_{\mathcal{F}}(\mathcal{G}),d_{\mathcal{G}}\big)\overset{\sim}{\longrightarrow}\big(\wedge^{\bullet}\mathfrak{h}^{*},d_{CE}\big).
		\]
		This implies in particular that $H^{1}_{\mathcal{F}}(\mathcal{G})\cong H^{1}(\mathfrak{h})$, so the proof is finished.
	\end{proof}

\end{document}
